\chapter{Detectores de semiconductor}

\section{Idea básica}
Un detector de semiconductor es una ``cámara de ionización de estado sólido'': volumen con campo
eléctrico donde la energía depositada crea \textbf{pares electrón--hueco} que derivan hacia los
electrodos e inducen una corriente en el circuito externo.

\begin{clave}[Ventajas frente a gas]
\begin{itemize}
\item \textbf{Energía por par muy baja}: 3.6 eV en Si, 2.96 eV en Ge (vs $\approx30$ eV en gas) $\Rightarrow$ muchos
portadores $\Rightarrow$ \textbf{excelente resolución en energía}.
\item \textbf{Densidad alta} (sólido): menor alcance para la misma energía, detectores compactos.
\item Campos altos y alta movilidad: \textbf{respuesta rápida} y buena eficiencia de colección (poco atrapamiento).
\item Resolución en posición dada por el tamaño de los electrodos (segmentación).
\end{itemize}
\end{clave}
\begin{dato}[gas / líquido / sólido]
Densidad: baja / moderada / alta. $Z$: bajo / moderado / moderado. Energía de ionización: moderada / moderada / baja.
Velocidad de señal: moderada / moderada / rápida.
\end{dato}

\section{Física del semiconductor}
\begin{itemize}
\item Si y Ge: 4 electrones de valencia, red tipo diamante. Gap: Si 1.12 eV, Ge 0.67 eV ($k_BT\approx25$--40 meV a temperatura ambiente).
\item \textbf{Dopado tipo n}: impureza pentavalente (P, As): donores, electrones mayoritarios.
\textbf{Tipo p}: impureza trivalente (B): aceptores, huecos mayoritarios.
\item Problema: en Si intrínseco $n_i\approx1.45\cdot10^{10}$ cm$^{-3}$ $\Rightarrow$ en un volumen típico hay $\sim4.5\cdot10^8$ portadores libres frente
a sólo $3.2\cdot10^4$ pares creados por una MIP. \textbf{Hay que eliminar los portadores libres: deplexionar} el detector
con una unión p--n polarizada en inversa.
\end{itemize}

\subsection{Unión p--n}
Por difusión, electrones del lado n pasan al p y huecos al revés, dejando carga espacial fija
(donores ionizados $+$ en n, aceptores $-$ en p). Se forma una \textbf{zona de deplexión} sin portadores
libres y una barrera de potencial que frena la difusión.
\begin{itemize}
\item \textbf{Polarización directa} ($+$ a p): baja la barrera, gran corriente.
\item \textbf{Polarización inversa} ($+$ a n): aumenta la barrera y la zona deplexionada; sólo circula una pequeña
\textbf{corriente de fuga} (portadores minoritarios), que crece con $T$, el espesor y la irradiación: es ruido.
\end{itemize}

\begin{clave}[Zona de deplexión y tensión de deplexión total]
Ecuación de Poisson (unión abrupta p$^+$--n):
\[ -\frac{\mathrm{d}^2\phi}{\mathrm{d}x^2}=\frac{q_0}{\varepsilon\varepsilon_0}N_{eff} \]
La densidad de carga es constante en la zona deplexionada, el campo crece linealmente (máximo en la unión) y el potencial es parabólico.
\[
w=\sqrt{\frac{2\varepsilon\varepsilon_0}{q_0|N_{eff}|}V}\ \propto\sqrt{V_B},\qquad
V_{dep}=\frac{q_0}{\varepsilon\varepsilon_0}\,|N_{eff}|\,d^2\quad\left(\text{exacto: }\tfrac{q_0}{2\varepsilon\varepsilon_0}|N_{eff}|d^2\right)
\]
\begin{itemize}
\item La zona deplexionada \textbf{crece con la tensión} ($w\propto\sqrt{V_B}$).
\item Los pares creados en la zona \emph{no} deplexionada (sin campo) no dan señal:
\textbf{colección completa de carga sólo si el detector está totalmente deplexionado ($V_B>V_{dep}$)}.
\item $V_{dep}\propto N_{eff}\,d^2$: material más puro (alta resistividad) o más fino $\Rightarrow$ menor tensión.
\end{itemize}
\end{clave}
\begin{dato}
Si tipo n de $\rho=2$ k$\Omega$cm ($N_D\approx2.2\cdot10^{12}$ cm$^{-3}$), 300 $\mu$m: $V_{dep}\approx150$ V, fuga $<1\,\mu$A.
\end{dato}

\section{Detectores de silicio}
\begin{clave}[Si: principalmente partículas cargadas y rayos X]
\begin{itemize}
\item Gap 1.12 eV, 3.6 eV/par: funciona a temperatura ambiente.
\item $\rho=2.33$ g/cm$^3$; MIP: $\dEdx\approx3.8$ MeV/cm $\approx108$ pares/$\mu$m.
\item Alta movilidad ($\mu_e=1450$, $\mu_h=450$ cm$^2$/Vs) $\Rightarrow$ colección $<10$ ns. Velocidad de deriva $v=\mu E$.
\item Muy puro, rígido (estructuras finas autoportantes), tecnología microelectrónica barata.
\item $Z=14$ bajo: poco eficiente para $\gamma$ de alta energía.
\end{itemize}
\end{clave}

\subsection{Silicios segmentados}
Segmentando el contacto p$^+$ en \textbf{tiras} (\emph{strips}) leídas por canales independientes se obtiene la posición
(resolución $\propto$ paso de las tiras, $\sigma\approx\text{paso}/\sqrt{12}$). Con tiras perpendiculares en ambas caras
(\textbf{DSSD}, \emph{double-sided strip detector}) se obtiene $(x,y)$. Espesor típico \textbf{300 $\mu$m} (150--500 $\mu$m).
Geometrías: cuadrados, anulares (tipo CD), en cajas alrededor del blanco; se combinan con Ge para medir partículas y $\gamma$ en coincidencia.

\subsection{Telescopios $\Delta E$--$E$ e identificación de partículas}
\begin{clave}[Telescopio $\Delta E$--$E$]
Un detector \textbf{muy delgado (5--50 $\mu$m)} mide la pérdida $\Delta E$ y otro \textbf{grueso (200--1500 $\mu$m)} detrás
para completamente la partícula ($E$ residual). \textbf{Sumando ambas señales se obtiene la energía total}.

Por \textbf{Bethe--Bloch}, en régimen no relativista
\[ \Delta E\propto\frac{Mz^2}{E}\quad\Rightarrow\quad \Delta E\cdot E\propto Mz^2 \]
cada especie (p, d, t, $^3$He, $\alpha$\dots) ocupa una \textbf{hipérbola distinta} en el diagrama $\Delta E$--$E$:
identificación de partículas. Seleccionando una banda se aísla un canal de reacción (p.ej.\ $^{163}$Dy($^3$He,$^4$He)$^{162}$Dy).
\end{clave}

\subsection{Relación señal/ruido}
La distribución de pérdida en capas finas (Landau) tiene cola a baja energía, ensanchada por el ruido. Fuentes de ruido
(ENC, carga equivalente de ruido): capacidad, corriente de fuga, ruido térmico de la resistencia de polarización.
Se exige un umbral por encima de la cola de ruido (alto $\Rightarrow$ pérdida de eficiencia; bajo $\Rightarrow$ ocupación por ruido).
Figura de mérito: $S/N>10$--15. El daño por radiación la degrada.

\section{Detectores de germanio (HPGe)}
\begin{clave}[HPGe: espectroscopía $\gamma$ y X de alta resolución]
\begin{itemize}
\item $Z=32$ (vs 14 del Si) y cristales grandes (coaxiales, cientos de cm$^3$): adecuado para $\gamma$.
\item Gap pequeño (0.67 eV) $\Rightarrow$ \textbf{hay que enfriar a 77 K (N$_2$ líquido)} o con refrigeración eléctrica para suprimir la
generación térmica de pares (corriente de fuga).
\item Alta pureza ($\sim10^{9}$--$10^{10}$ impurezas/cm$^3$, frente a $10^{12}$ en los antiguos Ge(Li), que necesitaban estar siempre fríos):
permite deplexionar varios cm con $V>1000$ V.
\item \textbf{Resolución $\approx0.15\%$ a 1.33 MeV} (FWHM 1.7--2.3 keV) frente a $\approx7$--8\% del NaI. Mejor cociente pico/Compton.
\item \textbf{Eficiencia intrínseca menor} que la de los centelleadores inorgánicos de $Z$ y $\rho$ altos.
\end{itemize}
\end{clave}
\begin{itemize}
\item Geometrías planar y coaxial; contactos de Li difundido (n$^+$, $\sim0.5$ mm, capa muerta gruesa) o B implantado (p$^+$, $\sim0.5\,\mu$m).
\item Para X y $\gamma$ de baja energía la ventana de entrada debe ser fina: contacto p$^+$ delante $\Rightarrow$ cristal tipo n
(además, el tipo n es más resistente al daño por neutrones). Ventana de Be para bajas energías, de Al para altas.
\item Cadena: detector $\to$ preamplificador $\to$ amplificador $\to$ ADC $\to$ ordenador. Constantes de conformado 2--6 $\mu$s.
\end{itemize}
\begin{dato}[Parámetros estándar a 1332 keV ($^{60}$Co)]
Eficiencia relativa (respecto a NaI 3''$\times$3'' a 25 cm, $\varepsilon=1.2\cdot10^{-3}$): 10--70\% (hoy $>100\%$).
Pico/Compton 30--60. FWHM 1.7--2.3 keV. Forma: FWTM/FWHM $\approx2.0$ (gaussiana 1.82).
\end{dato}

\begin{clave}[Detectores HPGe compuestos y segmentados]
\begin{itemize}
\item Para \textbf{mejorar el cociente pico/Compton} y \textbf{corregir el desplazamiento Doppler} (fuentes en vuelo) hay que conocer
dónde interacciona el $\gamma$: \textbf{tracking} $\Rightarrow$ cristales de Ge \textbf{segmentados} + análisis de forma de pulso.
\item Para aumentar la eficiencia se agrupan en \textbf{arrays} (clústeres con criostato común; detectores \emph{clover}).
\item Supresión Compton: escudos de BGO alrededor que vetan los sucesos en que el fotón escapa.
\end{itemize}
\end{clave}
\begin{dato}[Ejemplos]
Euroball, Exogam (clovers), Miniball (ISOLDE, cristales segmentados en 6), \textbf{AGATA} y \textbf{GRETINA/GRETA}
(tracking: 36 segmentos por cristal, 6 longitudinales $\times$ 6 transversales).
\end{dato}

\section{Otros semiconductores}
CdTe, CdZnTe (CZT), HgI$_2$: $Z$ alto y gap mayor $\Rightarrow$ \textbf{no requieren refrigeración} ($\varepsilon\approx4.4$ eV/par),
pero peor transporte de carga (huecos); por ahora para energías bajas. Diamante, GaAs, SiC: resistentes a la radiación.
SiPM (fotomultiplicadores de silicio): ver Tema 5.
